\section{Completion runs, scheduling and persistent rules}
\label{sec:p12-fair-completion}

A final critical-pair check proves correctness of a completed rewrite
system. An incremental completion procedure also needs to justify its
bookkeeping: earlier joining rules may be simplified or deleted before
the procedure stops. We distinguish Knuth--Bendix's restart procedure
\cite[\S\S5--6]{KnuthBendix1970} from Huet's labeled and marked
incremental procedure \cite[pp.~13--19]{Huet1981CompletionCorrectness}.
The latter supplies the fairness and limit results proved here. It is
a specified source algorithm, not a correctness assertion about every
possible interleaving of the primitive operations on
Sheet~\ref{sheet:term-rewriting}.

\begin{hbconvention}{def:p12-fc-common-order}{One ordering for the entire run}
All terms and rules have the ordinary first-order meaning of
Definition~\ref{def:p12-trs}. The signature and finite terms are
effectively represented. Fix one decidable reduction ordering \(>\),
well-founded and closed under substitutions and contexts, before the
run begins. Every generated rule must decrease this same ordering.
It may be partial. A nontrivial normalized equation whose sides are
incomparable causes failure in the oriented procedures below.
For the original 1970 presentation the ordering is
\(>_{\rm kb70}\); the proofs apply to any fixed ordering with the
stated properties. The input equation set \(E_0\) is finite.
\end{hbconvention}

\begin{hbmethod}{def:p12-fc-restart}{The restart and interreduction procedure}
First orient the input equations, discarding identical sides and
reporting failure if a required pair is incomparable. Interreduce the
finite rule set: whenever a rule has a side reducible by the other
rules, normalize both sides using those other rules. Delete the rule
if the two results agree; otherwise replace it by their greater-to-lesser
orientation, or report failure if they are incomparable.

After interreduction, compute every critical pair of the current rule
set, including both ordered overlap directions and self-overlaps.
Normalize each pair's two sides. If every resulting pair agrees, stop
with success. Otherwise orient a disagreeing pair, add that consequence
as a rule, interreduce again, and restart the complete critical-pair
scan. A fixed finite scan suffices for each current finite rule set;
no result from an earlier scan is reused without a new check.
These are the restart and rule-replacement operations of the 1970
source. Repeating interreduction until no replacement applies makes
its all-rules replacement instruction operationally definite; the
source does not prescribe the order of those replacements.
Its suggestions about trying another ordering or introducing
defined operators after failure are different runs or presentations,
not unrecorded changes inside this run.
\end{hbmethod}

\begin{hbtheorem}{thm:p12-fc-restart}{proved=annex:p12.fc.restart}{Correctness of successful restart completion}
Each interreduction phase terminates or reports failure. Every
successful restart output is terminating, confluent and equationally
equivalent to \(E_0\). The exhaustive final scan proves the required
critical-pair coverage; joinability is not an extra assumed postcondition.
This assertion does not say that repeated restarts eventually succeed.
\end{hbtheorem}

\begin{hbdefinition}{def:p12-fc-state}{Huet's labeled and marked states}
At iteration \(i\), a state consists of a finite set \(E_i\) of
pending equations, a finite set \(\mathcal R_i\) of records
\((k,\ell,r,m)\), and a counter \(c_i\). A record denotes the
rule \(k:\ell\to r\), whose mark \(m\) is either unmarked or
marked. Labels are unique natural numbers, allocated in increasing
order and never reused. Let \(R_i\) be the underlying rules after
forgetting labels and marks. A surviving label keeps its left-hand
side; its right-hand side may be normalized without changing its
label or mark. A changed left-hand side deletes that label.
Initially \(i=0\), \(E_i=E_0\),
\(\mathcal R_i=\varnothing\), and \(c_i=0\).
The following macrosteps define the subsequent iterations.
\end{hbdefinition}

\begin{hbmethod}{def:p12-fc-equation-phase}{Equation reduction and insertion}
While \(E_i\ne\varnothing\), select any \(s\approx t\in E_i\)
and compute \(R_i\)-normal forms \(s'\) and \(t'\).
If \(s'=t'\), remove the selected equation and leave the rules
unchanged. If \(s'\ne t'\) and neither side exceeds the other,
stop with failure. Otherwise let \(\ell>r\) be their orientation.
Put
\[
 K=\{k:(k,\ell_k,r_k,m_k)\in\mathcal R_i,
                   \ \ell_k\text{ is reducible by }\ell\to r\}.
\]
For every \(k\in K\), choose one such step
\(\ell_k\to_{\{\ell\to r\}}\ell'_k\). The next pending set is
\[
 E_{i+1}=(E_i\setminus\{s\approx t\})
                 \cup\{\ell'_k\approx r_k:k\in K\}.
\]
Allocate label \(c_i+1\) to the new, unmarked rule \(\ell\to r\).
Delete all records whose labels lie in \(K\). For each other old
record retain its left side, label and mark, replacing its right side
by any normal form under \(S_i=R_i\cup\{\ell\to r\}\).
These choices of right-side normal forms refer to this displayed
pre-insertion rule set plus the new rule. Finally increase the counter
and the iteration index. Ordinary set union may remove duplicate
pending equations.

The new left side is \(R_i\)-normal before insertion. This condition
controls which old rules may be collapsed and is essential to the
limit argument. No arbitrary removal of a rule is part of the procedure.
\end{hbmethod}

\begin{hbmethod}{def:p12-fc-critical-phase}{Incremental critical-pair computation}
When \(E_i=\varnothing\), stop with success if all rule records are
marked. Otherwise choose an unmarked label \(k\). Compute all
critical pairs between its current rule and every current rule whose
label \(j\leq k\), in both overlap directions, including \(j=k\).
Set the next pending equation set to those pairs, mark record \(k\),
leave all other records and the counter unchanged, and resume the
equation phase. Equal-sided pairs may be retained for trivial deletion.
The mark is retained if a later insertion only changes the right-hand
side. This is the source's incremental policy, whose correctness
requires the proof below; it is not a claim that the old and new
critical pairs are literally equal.
\end{hbmethod}

\begin{hbdefinition}{def:p12-fc-fairness}{Fair label selection and the two limits}
A selection is \emph{fair} when every allocated label is eventually
either deleted or selected for critical-pair computation. It is enough
always to select the least unmarked live label. Unqualified preference
for a short or promising rule is not this condition.

Consider a run that never reports failure, either infinite or extended
constantly after a successful halt. Let
\[
 R_* =\bigcup_i R_i.
\]
Let \(R_\infty\) contain the final versions of precisely the labels
that are never deleted after their creation. The lemmas below prove
that each such label's right side changes only finitely often, so this
definition is meaningful. Equivalently, these are the labeled rules
present unchanged from some iteration onward. \(R_*\) includes
transient versions; \(R_\infty\) need not be finite and is not an
effectively recognizable set merely by being called a limit.
\index{completion!fairness}\index{completion!persistent rules}
\end{hbdefinition}

\begin{hblemma}{thm:p12-fc-invariants}{proved=annex:p12.fc.invariants}{Termination of phases and source invariants}
In every nonfailure run:
\begin{enumerate}
\item Every rule of \(R_*\) decreases \(>\), and every \(R_i\)
      is interreduced: each right side is \(R_i\)-normal and each
      left side is normal under the other rules.
\item Every equation phase ends with an empty pending set in finitely
      many iterations. No fairness assumption on equation selection
      is needed.
\item Every pending equation \(s\approx t\) occurring at any
      iteration has a join \(s\to_{R_*}^*u\leftarrow_{R_*}^*t\).
\item Every generated equation and rule is an equational consequence
      of \(E_0\), and \(\equiv_{R_*}=\equiv_{E_0}\).
\item A label that is never deleted has only finitely many proper
      changes of its right side, ending in a rule of \(R_\infty\).
\end{enumerate}
\end{hblemma}

\begin{hbdefinition}{def:p12-fc-containment}{Pattern containment and strict generalization}
A term \(a\) \emph{contains} \(b\) when some subterm of \(a\)
is a substitution instance of \(b\). Write \(a\gg b\) when
\(a\) contains \(b\) but \(b\) does not contain \(a\).
This is a relation on patterns, not the reduction ordering \(>\).
Equal-size examples can be strict: \(f(x,x)\gg f(y,z)\) for
distinct \(y,z\), since generalization can separate two variables.
\end{hbdefinition}

\begin{hblemma}{thm:p12-fc-containment}{proved=annex:p12.fc.containment}{Well-founded deletion chains and persistent reducibility}
The relation \(\gg\) is well-founded. If insertion of \(\ell\to r\)
deletes an old label with left side \(a\), then \(a\gg\ell\).
Every term reducible by \(R_i\) remains reducible by every later
\(R_j\). In particular deletion is backed by a strictly more general
or structurally smaller matching pattern; it cannot simply discard
the old rule's applicability.
\end{hblemma}

\begin{hblemma}{thm:p12-fc-fair-coverage}{proved=annex:p12.fc.fair-coverage}{What fair marking supplies for a persistent peak}
Assume fair label selection. For each one-step persistent peak
\(a\to_{R_\infty}b_1\), \(a\to_{R_\infty}b_2\), there are
\(d_1,d_2,e\) such that
\[
 a\to_{R_*}d_i,\qquad d_i\to_{R_*}^*b_i,
 \qquad d_i\to_{R_*}^*e\quad(i=1,2).
\]
The join is established for appropriate historical right-side versions.
It is not an assumption that the final critical pair was already
joinable in \(R_\infty\).
\end{hblemma}

\begin{hbtheorem}{thm:p12-fc-limit}{proved=annex:p12.fc.limit}{Correctness of the fair persistent limit}
For every nonfailure run with fair label selection, \(R_\infty\)
is terminating, confluent and interreduced, and
\[
 \equiv_{R_\infty}=\equiv_{R_*}=\equiv_{E_0}.
\]
More strongly, whenever \(a\to_{R_*}^*b\), the terms \(a,b\)
have a common \(R_\infty\)-reduct; and whenever
\(a\to_{R_*}^*b_1,b_2\), the two endpoints have a common
\(R_\infty\)-reduct. Thus transient right-side versions do not
leave unaccounted equational consequences in the limit.
\end{hbtheorem}

\begin{hbtheorem}{thm:p12-fc-computation}{proved=annex:p12.fc.computation}{Finite success and semidecision during an infinite run}
A successful finite run returns a finite canonical presentation of
\(E_0\). Under the nonfailure and fairness hypotheses, an infinite
run has an infinite \(R_\infty\); if that limit is finite, the
procedure stops with success.

During a nonfailure fair run, progressively normalize any two query
terms with the successive current rule sets. Equality is accepted as
soon as the two current terms coincide. This semidecides their equality
modulo \(E_0\). After successful termination, comparing the two
final normal forms decides equality and inequality. An incomparable
equation, timeout or unfinished fair run is not an inequality certificate.
\end{hbtheorem}

\begin{hbcorollary}{thm:p12-fc-antichain}{proved=annex:p12.fc.antichain}{The pattern obstruction to an infinite fair run}
In an infinite nonfailure fair run, the left sides of
\(R_\infty\) form an infinite set of pairwise incomparable
patterns under containment: no distinct left side contains another.
Consequently, if the given term language has no infinite set of
pairwise containment-incomparable left-side patterns, every
nonfailure fair run terminates successfully. This is a property of
the fixed ordinary term language, not a completeness assertion for
rewriting modulo additional equations
\cite[p.~19]{Huet1981CompletionCorrectness}.
\end{hbcorollary}

\begin{hbtheorem}{thm:p12-fc-scheduler}{proved=annex:p12.fc.scheduler}{An explicit fair scheduler and a worked collapse run}
Selecting the least unmarked live label is fair. With input
\(\{f(a)\approx b,\ a\approx c\}\) and the reduction ordering
from Theorem~\ref{thm:p12-completion-worked}, Huet's procedure can
produce labels for \(f(a)\to b\), \(a\to c\), and then
\(f(c)\to b\), deleting the first label when the second is inserted.
The returned rule set is \(\{a\to c,f(c)\to b\}\), and the
marking scan establishes its complete critical-pair coverage.
\end{hbtheorem}

\begin{hbtheorem}{thm:p12-fc-associativity}{proved=annex:p12.fc.associativity}{The original associativity superposition test}
Take a finite signature consisting of one binary operation
\(\mathbin{\cdot}\) and at least one constant. The single rule
\[
 (x\cdot y)\cdot z\longrightarrow x\cdot(y\cdot z)
\]
is a canonical presentation of associativity. The original
\(>_{\rm kb70}\), with any strict total symbol precedence and
every symbol weight and the variable weight equal to one, orients
it. Its only nontrivial critical pair
has both sides reducing to \(x\cdot(y\cdot(z\cdot w))\).
Both specified procedures therefore succeed without adding any
rule. The canonical forms preserve the left-to-right sequence
of variable and constant leaves \cite[\S5, p.~276]{KnuthBendix1970}.
\end{hbtheorem}

\begin{hbtheorem}{thm:p12-fc-permutation-failure}{proved=annex:p12.fc.permutation-failure}{Commutativity cannot be oriented as one universal rule}
For a binary symbol \(f\) and distinct variables \(x,y\), no
substitution-stable strict ordering compares
\(f(x,y)\) and \(f(y,x)\) in either direction.
Thus the initial singleton equation
\(f(x,y)\approx f(y,x)\) causes an incomparable-equation
failure in either specified oriented procedure. This failure does
not decide that the equation is false, or that its equational
word problem is undecidable \cite[pp.~13--14]{Huet1981CompletionCorrectness}.
\end{hbtheorem}

\paragraph{Limits of the algorithmic claim.}
The fixed ordering matters. The singleton systems \(\{a\to b\}\)
and \(\{b\to a\}\) each terminate, but their union cycles and no
one strict ordering orients both. Stagewise termination alone therefore
cannot replace the common-order hypothesis.

Classical oriented completion stops at an incomparable nontrivial equation.
Unfailing completion retains or otherwise processes such equations with
additional inference rules. Ordered completion controls their use by
ordering conditions on suitable instances. Rewriting modulo a congruence
changes matching, unification and critical-pair formation. These are
distinct procedures requiring their own completeness arguments. In
particular, Huet's 1981 suggestions about extensions to permutative
or associative-commutative theories are not turned here into proved
claims about those extensions \cite[pp.~19--20]{Huet1981CompletionCorrectness}.
